$B$ is semisimple, so the kernel of a homomorphism $\theta : B \to A$ is a sum of Wedderburn blocks of $B$, and $\theta$ is injective on the complementary blocks; unitality ($\theta(1) = 1 \neq 0$) forces at least one block to survive. It suffices then that every block of $B$ has real dimension exceeding $\dim_\R A$. Rule 1, $A$ simple: $B \cong \Mat_2(A)$ is a single block of dimension $4 \dim A > \dim A$. Rule 1, $A = S \oplus S$ split: $\Mat_2(A) \cong \Mat_2(S) \oplus \Mat_2(S)$, blocks of dimension $4 \dim S = 2 \dim A > \dim A$. Rule 2 ($A$ split): $B \cong \Mat_2(\Cl^0)$ is simple (Theorem \claimref{transport.rule-2-ideal-fusion-sharp-form}) of dimension $2^{n+1} = 2\dim A > \dim A$. In every case no block embeds in $A$, contradiction.
